O presente relatório abordou, de forma prática e teórica, os principais métodos numéricos utilizados para a resolução de equações não lineares (zeros de funções) e sistemas de equações lineares. Através da implementação computacional e da análise dos resultados, foi possível compreender a fundo o comportamento, as vantagens e as limitações de cada técnica matemática, reforçando os conceitos vistos em sala de aula de Análise Numérica.

No que tange aos métodos para encontrar raízes de funções (Bisseção, Falsa Posição, Ponto Fixo, Newton-Raphson e Secante), ficou evidente a relação direta (trade-off) entre garantia de convergência e velocidade (taxa) de convergência. Enquanto métodos intervalares como a Bisseção garantem a raiz à custa de uma convergência linear e lenta, métodos abertos como Newton-Raphson e Secante mostraram-se formidáveis em eficiência (convergência quadrática e superlinear, respectivamente), mas sob a penalidade de requererem bons palpites iniciais e, por vezes, cálculo de derivadas. A análise gráfica e as tabelas com o histórico de iterações permitiram visualizar tais características, corroborando que não há um método "perfeito", mas sim o mais adequado para cada topologia de função.

Para os sistemas lineares, a investigação abraçou tanto os métodos diretos (Eliminação de Gauss e Fatoração LU) quanto os iterativos (Gauss-Jacobi e Gauss-Seidel). Os métodos diretos mostraram-se eficientes para sistemas menores e melhor condicionados, entregando a solução exata (a menos de erros de arredondamento) num número finito de passos. A decomposição LU destacou-se pela excelente reutilização dos cálculos quando se tem múltiplas resoluções com a mesma matriz dos coeficientes. Em contraponto, os métodos iterativos revelaram o seu valor principalmente ao lidarem com matrizes esparsas ou de grande porte, quando as condições de convergência (como o critério das linhas ou de Sassenfeld) são correspondidas, economizando memória computacional e proporcionando aproximações dentro de um erro tolerável em poucas iterações.

Ademais, foi elucidado o papel crucial do condicionamento de matrizes. Ao estudar a análise de condicionamento, as perturbações e os erros enfatizaram a importância de se utilizar pivoteamento parcial para amenizar instabilidades numéricas que podem arruinar toda a solução de um sistema aparentemente simples. 

Em suma, a experimentação desses algoritmos consolida a ideia de que a Análise Numérica é indispensável para as engenharias e ciências da computação atuais. A escolha do método correto em problemas reais requer conhecimento não só do embasamento teórico-matemático (taxas de convergência, estabilidade e condicionamento), mas também o compromisso com os recursos computacionais disponíveis (memória e tempo de execução). As experiências obtidas ao longo da resolução destas questões fornecem a base necessária para a aplicação inteligente da computação na solução de problemas matemáticos complexos aplicados.